\section{Study Design and Experimental Setup}

\subsection{Controlled Isaac Sim Setting}

The study uses a common Isaac Sim execution architecture across navigation and
tabletop manipulation \cite{nvidiaIsaacSimDocs}. The navigation family
contains language-conditioned mobile-goal tasks in indoor scenes. The
manipulation family contains tabletop tasks that require object identification,
grasping, transfer, and placement. The main comparison uses a shared easy
slice plus one harder slice per family, letting us test whether additional task
burden changes the qualitative ordering or mainly amplifies planner/tool
overhead.

All conditions share the same high-level handoff:

\begin{enumerate}
\item The planner receives the current task state and the exposed tool set.
\item The planner emits an action request.
\item Under R0, the request is dispatched as issued; under R1, the runtime
validates it against the active contract and allows one repair attempt on
failure.
\item The executor runs only requests that satisfy the active dispatch interface.
\end{enumerate}

The dispatch vocabulary is family-specific and explicit. Navigation exposes
state getters and \texttt{navigate\_to}; manipulation exposes state getters
and \texttt{scripted\_pick\_place\_step}. The study therefore asks whether the
planner can hand the executor a dispatchable request under a specified
interface, not whether it can merely describe the task in natural language.

\subsection{Contract and Runtime Variants}

The controlled matrix is summarized in Table~\ref{tab:experimental-design-summary}.
P0 uses an under-specified direct-action contract, allowing the planner to
emit a task-level next action without binding it to the executor's tool
namespace. P1 uses a typed tool-call contract, requiring the planner to
choose one exposed tool and provide dispatchable arguments. P2 keeps the same
typed tool-call contract and adds a brief planner-side self-check before
emission. Table~\ref{tab:contract-interface-examples} records one planner
emission for each contract format.

R0 is a bare-dispatch runtime: it does not perform pre-dispatch validation or
offer a repair path. R1 is a validate-and-retry runtime: it checks the
emitted request against the active contract and, on failure, returns the
validation error to the planner and allows one retry. Contract and runtime
therefore play distinct roles. Contract design can prevent invalid actions at
the source, while runtime validation can recover some contract violations that
still slip through.

\begin{table}[t]
\centering
\caption{Experimental design summary. The main factorial comparison covers a shared easy slice plus one harder slice per task family. Two compact ablations isolate action-interface and runtime-validation effects without adding new experiment axes.}
\label{tab:experimental-design-summary}
\footnotesize
\resizebox{\columnwidth}{!}{%
\begin{tabular}{llllrrl}
\hline
Block & Family & Slice & Varying axes & Runs/cell & Runs & Role \\
\hline
Main Factorial & Navigation & Easy & P0/P1/P2 x R0/R1 & 11 & 66 & Main Contract X Runtime Ordering \\
Main Factorial & Manipulation & Easy & P0/P1/P2 x R0/R1 & 3 & 18 & Main Contract X Runtime Ordering \\
Main Factorial & Navigation & Harder & P0/P1/P2 x R0/R1 & 4 & 24 & Harder-Task Robustness \\
Main Factorial & Manipulation & Harder & P0/P1/P2 x R0/R1 & 3 & 18 & Harder-Task Robustness \\
Action-Interface Ablation & Mixed & Easy Shared Subset & P0/P1/P2 at fixed R0 & 4 & 12 & Invalid Actions And P2 Overhead \\
Runtime-Only Ablation & Mixed & Recoverable Easy Probes & R0/R1 at fixed P1 & 4 & 8 & Recovery And Retry Value \\
\hline
\end{tabular}%
}
\end{table}


\subsection{Implementation Snapshot}

The study uses a repo-local deterministic planner backend implemented in
\texttt{src/isaacsim\_agent/planner/mock.py} rather than a hosted
foundation-model API. Given the current task state, exposed tools, and any
retry message, it emits one JSON action per planner call; there are therefore
no temperature, top-$p$, or rollout-sampling settings in this evaluation.

P0 realizes a direct-action JSON contract, P1 realizes typed tool JSON, and
P2 keeps the same typed tool JSON while adding a planner-side
\texttt{self\_check} string. The planner request already exposes the valid tool
names, descriptions, and required arguments through \texttt{available\_tools}.
R1 adds one more element: if parsing or validation fails, it appends the
literal error string to a short repair instruction and allows one retry. The
runtime parses the emitted JSON, validates the tool name against the
task-specific registry, enforces the exact required argument set, and then
checks task-specific payload semantics. For navigation,
\texttt{navigate\_to} must carry a valid \texttt{target\_pose} that matches the
configured goal pose. For manipulation,
\texttt{scripted\_pick\_place\_step} must carry the next scripted
\texttt{phase\_name}. The added P2 \texttt{self\_check} field is recorded in the
planner trace but does not replace runtime contract checking.

Navigation exposes \texttt{get\_robot\_state}, \texttt{get\_goal\_state}, and
\texttt{navigate\_to}. The first two are state queries; \texttt{navigate\_to}
dispatches one deterministic executor step toward the configured goal and
returns the applied target pose, progress, moved distance, and stuck-step
count. Manipulation exposes \texttt{get\_gripper\_state},
\texttt{get\_object\_state}, \texttt{get\_target\_state}, and
\texttt{scripted\_pick\_place\_step}. The first three are state queries;
\texttt{scripted\_pick\_place\_step} advances one phase of a fixed pick-and-place
script such as \texttt{move\_to\_pregrasp}, \texttt{descend\_to\_grasp}, or
\texttt{close\_gripper\_and\_attach}, and returns the applied phase plus the
updated observation and motion deltas. The executor therefore exposes
task-scoped, deterministic action primitives rather than raw controls or
free-form code execution.

\subsection{Metrics and Evaluation Protocol}

Success rate is the primary task outcome. An invalid action is an emitted
action or tool call that violates the active contract and therefore cannot be
dispatched as issued. A recovery is a run that contains an invalid action but
still ends successfully after runtime handling. Retries capture the extra
repair work used in R1. Planner calls and tool calls together define
planner/tool overhead, which is the paper's efficiency proxy.

The planner backend is deterministic by design, so the empirical unit is the
task instance under a specified condition rather than repeated stochastic
rollouts. The four main comparison cohorts cover 21 fixed task instances
(11 navigation easy, 3 manipulation easy, 4 navigation harder, and 3
manipulation harder), evaluated once under all six contract/runtime cells for
126 main-matrix runs. Two focused ablation cohorts add four task instances
each to isolate the action-interface and runtime effects, giving 29 task
instances and 146 total condition evaluations in the full fixed evaluation
set. We therefore report descriptive comparisons over this fixed evaluation
set rather than confidence intervals or significance tests across i.i.d.
stochastic replications.

The paper draws its quantitative comparisons from the main factorial results
plus targeted runtime-only, action-interface ablation, and harder-task
summaries. The combined summary provides the aggregate view, while the
slice-specific summaries support the ablations and the failure-case
interpretation. Within each slice, every compared condition is evaluated on the
same fixed tasks, which is why this task coverage still supports a controlled
systems comparison even without stochastic rollout replication.
